\chapter{Abstract}
% A brief summary of the study, including objectives, methods, key findings, and implications.

Wearable augmented reality (AR) devices, such as the Ray-Ban Meta glasses, allow users to take pictures, record videos from the wearer's perspective, and share them online with friends and family on social media. The glasses are also able to make phone calls, stream music, and have an AI assistant that allows users to perform these tasks hands-free. However, these capabilities also present unprecedented privacy concerns for users and bystanders. This thesis investigates the privacy implications of the Ray-Ban Meta glasses through static and dynamic analyses. The static analysis focuses on the device and its companion app, while the dynamic analysis examines network traffic and Bluetooth communications captured during the experiments.

The study reveals that, even though the captured traffic has strong encryption protocols protecting the transmitted content, metadata, such as packet sizes and traffic patterns, can still reveal application usage and sensitive information, including media transfers, certain command executions, and pairing events. Bluetooth communications further allow for device identification and potential user profiling. The exposed information may also include data related to bystanders, who are often unaware of the collection and, therefore, unable to provide informed consent. Although the companion app did not request excessive permissions nor include unnecessary third-party libraries, it continuously communicated with Meta's servers. These continuous transmissions raise concerns about potential long-term tracking, analytics, and possible unauthorized access to the data, which could similarly affect bystander privacy. While the encryption restricted a complete analysis of the traffic, partial decryption still revealed the existence of detailed user data, reinforcing the privacy concerns made in the metadata analysis.


